\subsection{The original integral and the signed boundary bound}
\label{sec:compact-model-signed-stokes}

\begin{lemma}\label{lem:compact-model-integral}
Let $Y$ be the original smooth complex projective variety, and let
$X$ be a supplied compact complex manifold of the same dimension,
whose projection from homogeneous coordinates has precisely the
coordinate-scaling orbits as fibres and admits local holomorphic
sections. Then $X\simeq Y$ biholomorphically. For an absolutely
integrable mixed curvature density, the original sum over inverse
projection branches equals its integral on $X$, with all complex
components included.

Moreover, an original vertical metric factor defined by a rational
section that is a regular unit away from one arithmetic prime has
zero curvature on $X$, including its complex boundary.
\end{lemma}

\begin{proof}
Local sections define a holomorphic bijection $X\to Y$. It is
biholomorphic by the injective holomorphic mapping theorem
\cite[Theorem~1.6.6]{compact-stokes-lebl}. Smoothness of $Y$ here is
the retained smooth-generic-fibre hypothesis.

A linear projection whose centre misses the Segre embedding of $Y$
is finite: its fibres are proper subvarieties of affine spaces,
hence finite \cite{compact-stokes-projection}. Outside its critical
values, infinity and the prescribed null boundary image, apply
change of variables to every inverse holomorphic branch. Different
homogeneous lifts change logarithmic metric potentials by
pluriharmonic functions. With $dd^c=(i/\pi)\partial\bar\partial$,
their mixed Hessian density is exactly the wedge of the curvature
forms. Applying this argument to the absolute value first justifies
the signed identity. Each projective point occurs once; any later
finite-cover multiplicities and stabilizer factors retain their
original values. In dimension zero this is a finite sum.

At every complex point the arithmetic prime is invertible. The
original unit relation supplies a local homogeneous quotient $G/H$
with $G,H$ nonzero and vertical factor $|G/H|$. Its logarithm is
locally the real part of a holomorphic logarithm, so its curvature
vanishes there. This argument applies at boundary points as well
\cite{compact-stokes-source}.
\end{proof}

Fix the original finite vertical cutoff and put $V$ for its product
of vertical factors. Let $h_0$ be the common reference denominator,
$\theta$ its curvature, and let the prescribed positive metric
$A'$ have Kähler curvature $\omega$. Suppose the tied potentials
$\varphi_t,\psi_{t'}$ satisfy the preceding boundary-energy lemmas,
with limits $\varphi,\psi$ and
$\psi-\varphi=U/2$. Write $T=\theta+dd^c\varphi\geq0$.
The actual smooth lines have restricted denominators
\[
 h_{\mu_t}=h_0e^{\varphi_t}V,\qquad
 h_{N_{p,q,t}}=h_{\mu_t}^{q}h_{A'}^p,\qquad
 h_{M_{p,q,t'}}=(h_0e^{\psi_{t'}}V)^qh_{A'}^p.
\]
Retain eventual smoothness of $\mu_t$ and, for each
$1\leq p\leq q$, of $M_{p,q,t'}$. The full intersection pairing
and all the displayed geometric data are supplied hypotheses.

\begin{proposition}[Signed arithmetic boundary comparison]
\label{prop:compact-model-signed-stokes}
Put $d(z)=\int_X(T+z\omega)^n$ and assume $d(0)>0$.
Suppose the original normalized height inequalities give
\[
 q^{-(n+1)}I(M_{p,q,t'}^{n+1})
       \geq(n+1)(a-(p/q)K_A)d(p/q)
\]
eventually for each $1\leq p\leq q$.
Then the original arithmetic limit exists and satisfies
\[
 \mathcal A:=\lim_{t\to0}I(\mu_t^{n+1})
 \geq(n+1)\left(a-\frac{\overline U}{2}\right)d(0),
 \qquad \overline U=\frac{\int_XU\,T^n}{d(0)}.
\]
\end{proposition}

\begin{proof}
Metric locality and multilinearity give one polynomial
\[
 F_t(z)=\sum_{j=0}^{n+1}\binom{n+1}{j}z^j
          I(\mu_t^{[n+1-j]},(A')^{[j]}),
 \qquad F_t(p/q)=q^{-(n+1)}I(N_{p,q,t}^{n+1}).
\]
The model, cutoff, vertical factors and metrics in this polynomial
are fixed independently of $p,q$. At each fixed rational
$z\in(0,1]$, exact metric variation and energy convergence give
a limit for $F_t(z)$. Values at $n+2$ distinct such points recover
all coefficients by interpolation. Thus $F_t\to F$ coefficientwise,
and $F(0)=\mathcal A$.

For fixed $z=p/q>0$, the bounded energy convergence gives a limit
$G(z)$ for the corresponding normalized bounded intersection.
No bounded line with $p=0$ is needed. The integral dictionary and
flatness of $V$ give the exact difference
\[
 G(z)-F(z)=(n+1)\bigl(E_z(\psi)-E_z(\varphi)\bigr).
\]
Here $E_z$ is unnormalized energy in the class of $\theta+z\omega$.
Near semipositivity supplies a smooth positive reference in that
class; $\theta$ itself need not be positive. The logarithmic metric
ratio contributes $q$ and the $n$ curvature factors contribute $q^n$,
exactly cancelling the original $q^{n+1}$.

Finite-energy concavity and $\psi-\varphi=U/2$ yield
\[
 E_z(\psi)-E_z(\varphi)
       \leq\frac12\int_XU\,(T+z\omega)^n.
\]
Together with the retained height inequality this gives
\[
 F(z)\geq(n+1)\left((a-zK_A)d(z)
                    -\frac12\int_XU\,(T+z\omega)^n\right).
\]
Full mixed mass identifies the geometric degree with $d(z)$.
Weighted mixed integrability makes the last integral a finite
polynomial in $z$. Letting positive rational $z$ tend to zero proves
the assertion with the original signed $U$. This also covers $n=0$.
\end{proof}

This proves the signed comparison for the given arithmetic and
regularization data. The construction of those data and the
arithmetic height and section estimates remain separate inputs.
